\textbf{Energía, ambiente y protección: decisión y capacidad conjunta.}
La infraestructura de continuidad se proyecta con dos caminos A/B de 400/230 V, 50 Hz, tres fases, neutro y conductor de protección, separados del resto del edificio. Cada camino comprende generación, transferencia, UPS, banco y distribución propios. T-11 conserva modelos, cantidades, variantes y condiciones de suministro; el Anexo A4.3 reúne las memorias de cálculo eléctricas, ambientales y de incendio. Esta separación permite evaluar la arquitectura y consultar cada partida sin repetir el inventario \parencite[RT-06.06--19, pp.~15--16]{pucv-transversal}.

La selección de referencia comprende dos UPS Eaton 93T60 de 60 kVA/60 kW y dos grupos FG Wilson P200-6 de 180 kVA/144 kW PRP, uno por camino. Cada UPS considera ochenta bloques CSB HRL12540W en dos cadenas independientes. El cálculo de selección desde la tabla de batería es $80(1.129)(0,85)(0,90)/1000=69,0948$ kW durante sesenta minutos en las condiciones declaradas de la memoria. No es una curva de autonomía del sistema: la aceptación debe comprobar carga, temperatura, pérdidas y degradación, manteniendo al menos treinta minutos a plena carga. La generación tiene objetivo de veinticuatro horas por camino, con estanque y reabastecimiento dimensionados separadamente \parencites{p7-eaton-93t60,p7-fg-p200-ficha,csbHRL12540WRA260131}[RT-06.07/08, p.~15]{pucv-transversal}.

\textbf{Resultado del balance integrado.} La memoria INT-40 del Anexo A4.3 conserva una base AC parcial de 75,108815 kW en régimen y 78,137615 kW en solape, antes de frío, bombeo primario y receptores restantes. Las diferencias respecto de 60 kW son $75,108815-60=15,108815$ kW y $78,137615-60=18,137615$ kW. Por ello la UPS de referencia no satisface todavía la combinación completa; su capacidad y autonomía no se declaran conformes. El cierre exige seleccionar y conciliar una sustitución o una distribución de cargas que cumpla cada estado, fases, kW/kVA, arranque y pérdida de camino. Las reservas DC se convierten a entrada AC una sola vez y no se suman nuevamente como una carga independiente.

El control ambiental ordinario reserva 356,7552 W DC, incluidos 4 W de desacoplamiento. En pérdida de una fuente, la superviviente debe sostener $136+80,7552+4=220,7552$ W: frente a 240 W nominales dispone de 19,2448 W de margen. A 65 grados C su capacidad de 210 W no alcanza esa demanda; acondicionamiento y arranque deben verificarse. Esta continuidad del mando no demuestra continuidad del frío, del circuito de hidrógeno ni de la alimentación de incendio. La matriz de funciones de INT-40 identifica esas fronteras.

\textbf{Climatización y protección independientes.} Se estudian dos circuitos de precisión para TI/UPS y dos fuentes de agua fría Carrier 30RB065-SS, con extracción y tratamiento separados por banco. La consigna de proyecto es 24 grados C y 45\,\% de humedad relativa, dentro de bandas de 23--25 grados C y 40--55\,\%, sin condensación. Cada fuente superviviente debe cubrir por sí sola la demanda autorizada. La demanda ideal de frío de dos rutas, 47,033856 kW, alcanza 48,473856 kW al incorporar una bomba secundaria; cuatro rutas a plena demanda alcanzan 94,067712 kW y exceden los 63 kW nominales de una fuente. Ese solape queda excluido; curvas corregidas, bombeo primario, estabilidad, tratamiento húmedo y transferencia permanecen por cerrar en el Anexo A4.3.

La protección mantiene detección, alarma local, extinción, portátiles y supervisión diferenciadas. El monitoreo no gobierna la liberación del agente. La aspiración ordinaria no se habilita en bancos de hidrógeno sin aptitud acreditada para la clasificación efectiva. La alimentación INC-42 limita el ramal fuente/detector a un metro; las variantes, caída, baterías, subtensión, retención y coordinación gas/fuego se desarrollan en A4.3. La pérdida WAN conserva alarma y protección locales, con entrega posterior de registros; no acredita aviso remoto instantáneo \parencite[RT-06.16--19, p.~16]{pucv-transversal}.

\textbf{Frontera energética y aceptación.} El PUE compara energía total del recinto y energía TI interior durante el mismo período; exportaciones a puestos/campo, nube y ensayos exteriores se contabilizan por separado. Los valores 4,74 y 4,78 corresponden a escenarios parciales y no se adoptan como PUE de la combinación completa. El balance anual debe incorporar frío, recarga, generación y variación de almacenamiento, sin sumar combustible químico como electricidad ni duplicar descarga y recarga. A4.3 conserva fórmulas, supuestos y escenarios para reconstruir ese cálculo \parencite[RT-06.11, p.~15; RT-15.04, p.~28]{pucv-transversal}.

Synaptix especifica, coordina, integra y administra; el CLIENTE adquiere hardware y ejecuta las obras asignadas. La recepción requiere planos, fichas de las variantes, protección/selectividad, ensayo de cada camino retirado, autonomía, transferencia, temperatura/humedad y alarmas. Las incompatibilidades indicadas bloquean la adquisición o habilitación del conjunto afectado; esta oferta no presenta esas pruebas como realizadas.
